\documentclass[11pt,a4paper,oneside]{book}
\usepackage{mathscommun}
\setniveau{Terminale}
\renewcommand{\ReferenceProgramme}{Spécialité (2026--2027 et 2027--2028) et mathématiques expertes}
\renewcommand{\EditionDocument}{Édition du 24 octobre 2026}
\hypersetup{pdftitle={Terminale - Cours et démonstrations}}
\begin{document}
\titrecollection{Cours et démonstrations}{Spécialité, mathématiques expertes et fondements pour le supérieur.}{Un parcours démontré : maîtriser les définitions, choisir les méthodes, contrôler les hypothèses et préparer la licence avant les spécialisations de master.}
\frontmatter
\chapter*{Mode d'emploi}
Ce volume suit le parcours de terminale générale avec spécialité mathématiques et option mathématiques expertes. Il ne constitue pas un cours des voies technologiques ni de l'option mathématiques complémentaires. La table de couverture du guide distingue les textes applicables en 2026--2027 et à partir de 2027--2028.

Les chapitres 1 à 24 forment le parcours de spécialité, avec les sections signalées en approfondissement, notamment les sphères et les barycentres du chapitre 8. Les chapitres 25 à 35 couvrent l'option expertes. Le chapitre 36 est une passerelle : il donne les fondements de plusieurs preuves et ne prétend pas remplacer une licence.

Les rappels de première sont repris au point d'usage. Pour une remise à niveau plus longue, revenir aux chapitres correspondants du volume de première : raisonnement (1), suites (4--7), dérivation (8--11), exponentielle (12), trigonométrie (14--15), produit scalaire (16), probabilités (19--23).

Lire une première fois pour comprendre les objets ; refaire ensuite les preuves sans support ; enfin chercher les exercices et les synthèses avant de consulter les corrigés. Chaque chapitre possède six exercices numérotés de la même façon dans les deux recueils. Les indications de difficulté distinguent les tâches de base, de raisonnement et d'approfondissement.

Les fondements admis sont recensés dans le guide. Les quantificateurs, la complétude et les notations abstraites supplémentaires sont explicités pour préparer le supérieur ; leur maîtrise complète dépasse les exigences scolaires.
\tableofcontents
\mainmatter
\subfile{chapitres/01_langage.tex}
\subfile{chapitres/02_recurrence.tex}
\subfile{chapitres/03_denombrement.tex}
\subfile{chapitres/04_binome.tex}
\subfile{chapitres/05_vecteurs.tex}
\subfile{chapitres/06_orthogonalite.tex}
\subfile{chapitres/07_equations_espace.tex}
\subfile{chapitres/08_geometrie_complements.tex}
\subfile{chapitres/09_limites_suites.tex}
\subfile{chapitres/10_operations_limites.tex}
\subfile{chapitres/11_convergence.tex}
\subfile{chapitres/12_limites_fonctions.tex}
\subfile{chapitres/13_continuite.tex}
\subfile{chapitres/14_derivation.tex}
\subfile{chapitres/15_convexite.tex}
\subfile{chapitres/16_logarithme.tex}
\subfile{chapitres/17_trigonometrie.tex}
\subfile{chapitres/18_primitives_ed.tex}
\subfile{chapitres/19_integrale.tex}
\subfile{chapitres/20_calcul_integral.tex}
\subfile{chapitres/21_binomiale.tex}
\subfile{chapitres/22_moments.tex}
\subfile{chapitres/23_concentration.tex}
\subfile{chapitres/24_algorithmes.tex}
\subfile{chapitres/25_complexes.tex}
\subfile{chapitres/26_complexes_trigo.tex}
\subfile{chapitres/27_polynomes.tex}
\subfile{chapitres/28_geometrie_complexe.tex}
\subfile{chapitres/29_division.tex}
\subfile{chapitres/30_bezout.tex}
\subfile{chapitres/31_premiers.tex}
\subfile{chapitres/32_matrices.tex}
\subfile{chapitres/33_puissances.tex}
\subfile{chapitres/34_graphes.tex}
\subfile{chapitres/35_markov.tex}
\subfile{chapitres/36_fondements.tex}
\end{document}
